\documentclass[openany,fontset=fandol]{ctexbook}
\usepackage{amsmath,amssymb}
\usepackage[hidelinks]{hyperref}
\newcommand{\note}[1]{\footnote{译注：#1}}
\title{代数几何基础第五卷预备笔记\\第12节：一般线性截面的推广说明}
\author{据用户上传英译编辑稿译出}
\date{2026年11月16日}
\begin{document}
\maketitle
\chapter*{范围与校订状态}
本批连续译出底本第44页“The Proposition 2.1”起，跨第45页，止于“Analogous change in 4.4 b)”。包括第2节各项、3.1—3.3及第4节有关条目的推广说明。下一处为第45页“[Le laius 5.1 Fr or Latin 67] goes as such”。

本段是原稿的编写提纲和推广指示，不是逐项重写的完整定理及证明。正文保留这一性质，不将“照旧成立”等作者说明改称本批已独立证明的结论。点的维数记号按前文澄清，无法辨认的不等式原样标为待核。此前8.8、9.5等缺文未解决；本批未合入原合订本。
\chapter*{\S12 将上述结果推广到线性截面（续）}
\section*{第2节的推广}
命题2.1可以不作改动地移用到这里〔原文法语：se transpose；英译在“translates”后留有问号〕。在2.2中，应当这样表述：充要条件是，对每个 $x\in Z$，都有
\[
\dim\overline{\{x\}}\leq m-1.
\]
\note{底本简写为“$\dim x\leq m-1$”。前文2.2讨论 $P$ 中子集在一般超平面中的逆像何时为空，其 $m=1$ 情形要求每一点都闭；证明也明确取点的闭包。因此这里把点的维数解释为其在 $P$ 中闭包的维数，而不是局部环维数或余维。对于域上有限型射影空间，这也等于该点剩余域在底域上的超越次数。此为记号澄清，不是在此补造全部推广证明。}
在证明中，例如可以把余维 $m$ 的一般线性子簇看成 $m$ 个独立的一般超平面的交，从而归约到2.6。

〔原稿插语〕让Dieudonné自己设法处理吧。〔原文：Dieudonné demerdetur，标作拉丁语；原英译者又问：这会不会其实是稍显粗俗的法语？〕\note{保留这句编写者之间的插语及原英译者对语种、语气的疑问；它不构成数学论证步骤。}

从编写的角度说，如果从一开始就取一般的 $m$（我觉得这样更好），那么似乎最好同时证明2.6；其中当然应把
\[
\dim X-1\quad\text{替换为}\quad\dim X-m,
\]
并约定公式中维数小于0表示所讨论的集合为空。\note{这是一条对原命题2.6的编写指示，必须连同该命题原有适用条件以及本节相应的非空性条件阅读。不能抽离条件后声称任意 $X$ 的任意余维 $m$ 截面都具有维数 $\dim X-m$。底本第8页2.6本身还有标为字迹不清的附语，本批没有将其恢复。}

推论2.3中，把“有限”替换为“维数不超过 $m-1$”，即可得到相应表述。推论2.4作类似处理。2.5也一样：把
\[
\dim f(X_i)>0\quad\text{替换为}\quad\dim f(X_i)\geq m;
\]
2.7中也作同样的替换。\note{2.7原文实际用 $Z_i$ 表示关联素循环，条件写为 $\dim f(Z_i)>0$；“同样的替换”是把这里的下界改为 $m$，不是将 $Z_i$ 改名为 $X_i$。本段各“有限”替换均针对相应原条目中的像集或例外集，仍须保留原条目的构造性、有限型等假设。}

命题2.8照原样仍然成立；在2.9中，将“有限”替换为“维数不超过 $m-1$”。2.10、2.11也作同样处理。

陈述2.12照原样仍然成立，证明基本不变（也参见下文关于第8节的说明）。在2.13中，把“有限”替换为“维数不超过 $m-1$”。2.14照原样成立；2.15则作上述“有限”的替换。2.16照原样有效；2.17也把“有限”替换为“维数不超过 $m-1$”。2.18照旧。\note{本段逐项保留原稿对2.8—2.18的指示，并未将各条的假设删去或重新统一。例如2.12原文有 $f$ 无分歧的假设。“证明基本不变”是原稿判断，不是本批对全部推广证明完成独立验证的声明。}

\section*{第3节与第4节的推广}
对于3.1，可以在假设
\[
\dim f(X)\geq m+1
\]
的条件下，对任意 $m$ 作出相应陈述。不过，我建议主要陈述仍保留在超曲面情形，并把一般情形作为推论或附记给出（用通常的办法，取独立的一般超平面，就能立即导出它）。

至少，把引理3.1.1的推广版本明确陈述出来会很有意思……对于3.2，条件应读作
\[
\dim f(X)\geq m+1.
\]
在3.3中，把 $\dim f(X_i)\geq2$ 替换为 $\dim f(X_i)\geq m+1$；并在 $G$ 的定义中，把 $\dim f(Z)=0$ 替换为 $\dim f(Z)\leq m-1$。\note{这里保留原稿的省略号及编写建议，没有凭引理编号补造3.1.1的推广陈述。此处 $G$ 是3.3中定义的对象，不应与本节开篇的泛商模 $G$ 混同。}

\enlargethispage{2\baselineskip}
第4节的一般讨论可以照原样用于任意 $m$。4.2和4.3也一样，只须在b）的(v)、(vi)中将维数条件作相应修改。〔底本在此写成：$\dim f(X)\geq m-1$，或者 $m+1$；字迹不清，请询问A.G.。〕4.4的b）作类似修改。\note{底本明确留下“$m-1$ or $m+1$ (illegible, ask AG)”的两种读法。本批未取得可消除歧义的原手稿，因此不选择其中一个作为正确正文条件，也不凭邻近条目的 $m+1$ 猜补；这一项仍待核。}
\end{document}
